\section{The Contrast-Concentration Theorem}

\begin{theorem}[Thalean quadratic contrast-concentration theorem]
Let
\[
\widehat{Q}
=
\frac{1}{98}
\left(A^3+2A^2+2I\right)
\]
on
\[
\mathbb R^{15},
\]
where \(A\) is the adjacency matrix of
\(\Gfifteen\cong L(\Petersen)\).

For every \(x\in\mathbb R^{15}\) and every integer \(k\ge0\):

\begin{enumerate}[label=(\roman*)]
\item the common component is fixed,
\[
\Pzero\widehat{Q}^{\,k}x
=
\Pzero x;
\]

\item the centered component satisfies
\[
\Rrel(\widehat{Q}^{\,k}x)
\le
\left(\frac{9}{49}\right)^k
\Rrel(x);
\]

\item the iterates converge to the common-mode projection,
\[
\lim_{k\to\infty}
\widehat{Q}^{\,k}x
=
\Pzero x
=
\overline{x}\,\one_{15};
\]

\item if \(\Ccommon(x)\neq0\), then
\[
\rho(\widehat{Q}^{\,k}x)
\le
\left(\frac{9}{49}\right)^k
\rho(x);
\]

\item if \(x\neq0\) and \(\Ccommon(x)\neq0\), then
\[
\zetaRel(\widehat{Q}^{\,k}x)
\le
\left(\frac{9}{49}\right)^{2k}
\rho(x)^2.
\]
\end{enumerate}
\end{theorem}

\begin{proof}
The matrix \(A\) is real symmetric, so its eigenspaces form an
orthogonal decomposition.  Since
\[
Q=A^3+2A^2+2I,
\]
the matrices \(A\), \(Q\), and \(\widehat{Q}\) have the same
eigenspaces.

On the all-ones eigenspace \(E_4\), the operator
\(\widehat{Q}\) has eigenvalue \(1\).  Therefore
\[
\Pzero\widehat{Q}^{\,k}x
=
\Pzero x.
\]

On the centered subspace
\[
\one_{15}^{\perp}
=
E_2\oplus E_{-1}\oplus E_{-2},
\]
the eigenvalues are
\[
\frac{9}{49},
\qquad
\frac{3}{98},
\qquad
\frac{1}{49}.
\]
All are nonnegative, and their maximum is \(9/49\).  Hence the
Euclidean operator norm of
\(\widehat{Q}^{\,k}\) restricted to the centered subspace is
\[
\left(\frac{9}{49}\right)^k.
\]
Thus
\[
\begin{aligned}
\Rrel(\widehat{Q}^{\,k}x)
&=
\left\|
\Pperp\widehat{Q}^{\,k}x
\right\|_2\\
&=
\left\|
\widehat{Q}^{\,k}\Pperp x
\right\|_2\\
&\le
\left(\frac{9}{49}\right)^k
\|\Pperp x\|_2.
\end{aligned}
\]

The centered bound tends to zero, while the common component remains
fixed.  Therefore
\[
\widehat{Q}^{\,k}x
=
\Pzero x+\Pperp\widehat{Q}^{\,k}x
\longrightarrow
\Pzero x.
\]

When \(\Ccommon(x)\neq0\), common-mode invariance gives
\[
\Ccommon(\widehat{Q}^{\,k}x)
=
\Ccommon(x),
\]
so division yields the bound for \(\rho\).

Finally,
\[
\zetaRel(y)
=
\frac{\Rrel(y)^2}
{\Ccommon(y)^2+\Rrel(y)^2}
\le
\frac{\Rrel(y)^2}{\Ccommon(y)^2}.
\]
Applying the preceding two bounds with
\(y=\widehat{Q}^{\,k}x\) proves the final inequality.
\end{proof}

\subsection{Exact modal formula}

Write
\[
x
=
x_4+x_2+x_{-1}+x_{-2},
\qquad
x_\lambda\in E_\lambda.
\]
Then
\[
\widehat{Q}^{\,k}x
=
x_4
+
\left(\frac{9}{49}\right)^k x_2
+
\left(\frac{3}{98}\right)^k x_{-1}
+
\left(\frac{1}{49}\right)^k x_{-2}.
\]

Consequently,
\[
\Rrel(\widehat{Q}^{\,k}x)^2
=
\left(\frac{9}{49}\right)^{2k}\|x_2\|_2^2
+
\left(\frac{3}{98}\right)^{2k}\|x_{-1}\|_2^2
+
\left(\frac{1}{49}\right)^{2k}\|x_{-2}\|_2^2.
\]

The bound in the theorem is sharp whenever the centered component
lies entirely in \(E_2\).

\begin{corollary}[Operator convergence]
In Euclidean operator norm,
\[
\widehat{Q}^{\,k}
\longrightarrow
\Pzero,
\]
and
\[
\left\|
\widehat{Q}^{\,k}-\Pzero
\right\|_{2\to2}
=
\left(\frac{9}{49}\right)^k.
\]
\end{corollary}

\begin{proof}
The difference vanishes on \(E_4\) and acts on the centered
eigenspaces with eigenvalues
\[
\left(\frac{9}{49}\right)^k,\quad
\left(\frac{3}{98}\right)^k,\quad
\left(\frac{1}{49}\right)^k.
\]
The largest absolute value is
\((9/49)^k\).
\end{proof}

\begin{remark}[Exactly centered input]
If
\[
\Pzero x=0,
\]
then every iterate remains centered:
\[
\Pzero\widehat{Q}^{\,k}x=0.
\]
In that case,
\[
\widehat{Q}^{\,k}x\longrightarrow0.
\]
The operator does not create a common component that was absent from
the input.
\end{remark}

\begin{remark}[Raw and normalized aggregation]
For the unnormalized quadratic,
\[
\Ccommon(Q^kx)
=
98^k\Ccommon(x)
\]
and
\[
\Rrel(Q^kx)
\le
18^k\Rrel(x).
\]
Thus the same ratio
\[
\frac{18}{98}
=
\frac{9}{49}
\]
governs relative contrast concentration under raw \(Q\)-iteration.

The normalized operator is preferable for the theorem because it
fixes the aggregate scale and turns the relative statement into an
absolute centered contraction.
\end{remark}

\subsection{Interpretive boundary}


The theorem proves exact algebraic convergence under repeated
application of a canonically normalized incidence operator.

It does not assert that native Thalean loading repeatedly applies
\(\widehat Q\), and it does not identify \(\widehat Q\) with physical
dynamics.

The native theorem now supplies a different exact operation:
\[
8\longrightarrow2\longrightarrow1
\]
for lawful completion-family contraction.

Under the canonical completion-occupancy register, native loading
strictly increases Paper 14 contrast on every strict loading
extension.

Therefore the former candidate identification
\[
\text{native loading}
\longrightarrow
\widehat Q\text{-type concentration}
\]
is rejected.

The correct relationship is
\[
\text{native selection}
\neq
\text{quadratic concentration},
\]
with both operations joined through the certified finite
transport-response algebra.
